\chapter{Game Theoretic Algorithms}
Game Theory is an \emph{idealisation} of real life games where under different axioms of the game and 
assumptions about the players' abilities we seek to find their optimal choices. 

\section{Definitions}
In a game each player has some \emph{actions} which they can play in the current \emph{state} when it is their turn. 
A state is a complete description of the current game situation. A \emph{history} is the sequence of 
states until (and including) the current state. \\

\noindent A \emph{strategy} for a player is a function that maps every possible state 
to an action which is legal to play. Simply stated, a strategy is a plan for what to play 
in what situation. A strategy may be as simple as "pick any random legal move" or as complicated 
as an entire mapping of every state to a carefully chosen legal 
move. A strategy profile is the collection of all strategies for all players. \\

\noindent A \emph{rational} player is a player who has the "intelligence" to find and choose a strategy that 
maximises their own utility. \\

\noindent A fact is \emph{common knowledge} in a set of players means that every player knows the 
fact, and every player knows that every player knows the fact and so on till infinity. This is an 
interesting definition but for our purposes the layman meaning will be sufficient.

\section{Nash Equilibria}
A Nash equilibrium is a situation where deviation from 
the chosen strategy by a player (keeping the strategies of the other players fixed) does not 
strictly increase the utility of the player. Such a deviation\footnote{In some sense, 
a Nash equilibrium is like the equilibrium of a body in physics where the energy is a function of multiple variables. To 
check equilibrium, a valid way would be to keep all variables except one fixed and analyse with 
respect to variation of that variable.} is called a \emph{unilateral deviation}. \\

We will use the idea of \emph{Nash Equilibrium} to predict the optimal action of a player under the following
assumptions (which will be referred to as the \emph{standard rationality assumptions}):

\begin{itemize}
\item All players are rational.
\item The fact that all players are rational is common knowledge among all players.
\end{itemize}

In case of multiple Nash equilibria, technically all of them are valid answers to the question -- "what 
will happen in the game under the standard rationality assumptions?"

\subsection{Kingdom's Dilemma}

The game is very simple with only one state and two actions for both players. The usual example is 
a symmetric game set between two kingdoms. Two kingdoms must independently choose between 
PEACE and WAR. They are not allowed to communicate or offer anything in 
exchange for cooperation or surrender - it is as if the only purpose of these two players is to 
maximise their utility and they care about nothing else. Let $S_A$ and $S_B$ be the strategies of 
Kingdom A and Kingdom B respectively.


\[
\begin{tabular}{|c|c|c|}
\hline
\diagbox{$S_B$}{$S_A$}
& \text{PEACE}
& \text{WAR}
\\
\hline
\text{PEACE}
& (4,4)
& (0,5)
\\
\hline
\text{WAR}
& (5,0)
& (1,1) \\
\hline
\end{tabular}
\]


\noindent Each entry in the above table is of the form $(U_A,U_B)$, where $U_A$ and $U_B$ denote 
the utilities of Kingdom A and Kingdom B respectively. Usually this table is a strategy-utility matrix but here the strategy is simply 
to choose an action so we wrote it as a action-utility matrix. This matrix is in general 
called the \emph{payoff matrix}.\\

\noindent The first coordinate denotes the utility of Kingdom A and the second denotes the utility 
of Kingdom B for this set of actions. So $(S_A,S_B)=(PEACE,PEACE)$ would give both $4$ utility, 
$(S_A,S_B)=(PEACE,WAR)$ would give A utility $0$ and B utility $5$, $(S_A,S_B)=(WAR,PEACE)$ would give 
A utility $5$ and B utility $0$ and $(S_A,S_B)=(WAR,WAR)$ would give both $1$ utility. \\

\noindent Looking at the table, we observe that it would of course be preferable for both to 
choose PEACE since that gives both of them $5$ utility. But a deviation from this would be 
beneficial -- for example, if Kingdom B chose WAR instead, they would gain $1$ extra utility. So 
clearly the strategy profile $(S_A,S_B)=(PEACE,PEACE)$ is not a Nash equilibrium. Looking 
at each of the four strategy profiles, it is clear that $(S_A,S_B)=(WAR,WAR)$ is a Nash equilibrium -- 
since any kingdom's change of strategy to PEACE would not increase their utility 
(in fact in this case it strictly decrease their utility). \\

\noindent This analysis reveals that a profile with better utilities for all players 
may not be the rational outcome! (PEACE,PEACE) was better for both players but was not an 
equilibrium$\dots$\\

\noindent A side point to be noted here is that a decrease of the utility of Kingdom A does not matter 
to Kingdom B. If B had a choice to go from $(U_A,U_B)=(3,5)$ to $(U_A,U_B)=(0,4)$ he would not 
take it since it decreases his own utility -- the common idea that another's loss is our gain 
is not true in general. The class of two-player games where this is true is called a \emph{constant-sum game}, 
because the sum of utilities of both players is constant for any strategy profile 
so maximisation of own utility automatically means minimisation of the opponent's. A common case of 
constant-sum game is the zero-sum game.

\subsection{Penalty Kick}
The shooter and the keeper are playing a game of penalty kick. The shooter is allowed to kick 
the ball in two sides -- left(L) and right(R). The keeper is allowed to choose between two actions -- jump 
to left(L) or jump to right (R). If both choose same side, the keeper wins. If they choose different 
sides, the shooter wins.\\

\noindent Let A be the keeper and B be the shooter. The payoff matrix is as follows:

\[
\begin{tabular}{|c|c|c|}
\hline
\diagbox {$S_B$}{$S_A$} & \text{L} & \text{R} \\
\hline
\text{L} & (-1,1) & (1,-1) \\
\hline
\text{R} & (1,-1) & (-1,1) \\
\hline
\end{tabular}
\]


Upon careful inspection, we notice that no square of the payoff matrix is a Nash equilibrium. This is 
because there exists an obvious deviation for the losing player that would 
strictly increase their utility -- just switch action to the other.\\

This should make us uncomfortable -- our planning stage would never end because there would be no 
settled conclusion of what strategy to have. So would a rational player just not play the game? Or would 
there be no point in choosing a strategy, just choose any one because does it really matter? This 
question itself suggests the introduction of randomisation as part of the strategy.

\section{Mixed Strategies}
A pure strategy is a deterministic map of states to actions. A mixed strategy is a mapping from states 
to a probability distribution over actions -- basically a particular probability of taking any action.
The failure of pure strategies suggests that perhaps randomisation would settle the game 
into equilibrium. In mixed strategies, a rational player seeks to maximise the \emph{expected} utility 
instead of the usual maximisation of utility.\\

The idea of a mixed strategy should feel familiar -- a deterministic algorithm can be 
exploited by an adversarial input and randomisation in the algorithm can help us deal with 
such adversaries, like in \texttt{QuickSort}.

\section{Mixed Strategy Nash Equilibria}
The Nash equilibrium we saw earlier is called a Pure Strategy Nash equilibrium. But with mixed strategies, 
the unilateral deviation is a player changing the probability distribution over actions. Note that 
the distribution can take any real values as long as the sum of probabilities is 1.

Let us consider the penalty kick game again. 

\[
\begin{tabular}{|c|c|c|}
\hline
\diagbox {$S_B$}{$S_A$} & \text{L} & \text{R} \\
\hline
\text{L} & (-1,1) & (1,-1) \\
\hline
\text{R} & (1,-1) & (-1,1) \\
\hline
\end{tabular}
\]

Let $S_A = (p:L \;,\; 1-p:R)$ be the distribution over L and R for player A. Similarly $S_B = (q:L\;,\;1-q:R)$ 
be the distribution over L and R for player B. The expected utility of player A can be written in terms 
of $p$ and the conditional expected utilities.

\[
\mathbb E[U_A] = p \cdot (E[U_A | S_A = L]) + (1-p) \cdot E[U_A | S_A = R]
\]
\[
\mathbb E[U_A | S_A = L] = q(E[U_A | S_A = L, S_B = L]) + (1-q)E[U_A | S_A = L, S_B = R] = q(-1) + (1-q)(1) = 1-2q
\]
\[
\mathbb E[U_A | S_A = R] = q(E[U_A | S_A = R, S_B = L]) + (1-q)E[U_A | S_A = R, S_B = R] = q(1) + (1-q)(-1) = 2q-1
\]

\[
\implies
\mathbb E[U_A] = p(1-2q) + (1-p)(2q-1) = (2p-1)(2q-1)
\]

Now to ensure unilateral deviation does not increase the expected utility of either player, we must 
have partial derivatives of the expected utility of player A with respect to $p$ and $q$ equal to zero.

\[
\frac{\partial \mathbb E[U_A]}{\partial p} = 0
\implies 
2(2q-1)=0 
\implies 
q = \frac{1}{2}
\]
\[
\frac{\partial \mathbb E[U_A]}{\partial q} = 0
\implies
2(2p-1)=0
\implies
p = \frac{1}{2}
\]

So both actions are chosen with equal probability for both players. This is not surprising 
because of the symmetries of the payoff matrix.\\

\noindent Now you may have noticed there was an easier way to find $E[U_A]$ by using the normal definition of 
expected value as a weighted average of the four cases \{LL,LR,RL,RR\}. But this way illustrates the fact that 
the partial derivative with respect to $p$ gave the value of $q$ and vice versa. This is 
not a coincidence, in fact it is rather obvious from the equation we saw earlier.

\[
\mathbb E[U_A] = p \cdot (E[U_A | S_A = L]) + (1-p)\cdot E[U_A | S_A = R]
\]

Here the conditional probabilities do not depend on $p$ so upon differentiation we get that both 
conditional expectations must be equal which upon solving yields the value of $q$. Applying the same 
argument on the corresponding equation written using $q$ would yield the value of $p$.

\[
\mathbb E[U_A] = q \cdot (E[U_A | S_B = L]) + (1-q)\cdot E[U_A | S_B = R]
\]

\section{Indifference Theorem}
The generalisation of the observation in the previous section is the \emph{indifference theorem}, where 
for a general set of probabilities $p_1,p_2,\dots,p_n$ (whose sum must obviously be 1) corresponding 
to $n$ strategies of player A, a Mixed Strategy Nash Equilibrium implies all the 
conditional expected utilities are equal. This makes player A indifferent to which strategy 
he chooses since all give same expected utility. If one of the conditional expected utilities is higher than the others, then 
player A will simply choose that strategy with probability 1. This may be viewed as B being 
forced to make all choices for A equal. Of course, the same thing applies with A being forced 
to make all choices for B equal. 

\[
\mathbb E[U_A] = \sum_{i=1}^n p_i \cdot E[U_A | S_A = i] \text { and MSNE }
\implies 
E[U_A | S_A = 1] = E[U_A | S_A = 2] = \cdots = E[U_A | S_A = n]
\]

An important question is -- when can mixed strategy Nash equilibria exist? When can both 
players enforce indifference for each other (or in general, multiple players)?

\section{The Minimax Principle}
We state without proof that for every finite two-player zero-sum game, there 
exists a mixed strategy Nash equilibrium. A pure strategy Nash equilibrium is 
considered to be a special case of the general mixed strategy Nash equilibrium.\\

\section*{Philosophical Note}
The analysis leading to Nash equilibrium assumes that when a player considers 
changing strategy, the strategies of all other players remain fixed. Some philosophers have argued 
that in perfectly symmetric situations between identical rational agents, this may not be the 
correct thing to check. If both agents reason in exactly the same way, they will expect both their 
planning stages to reach identical conclusions. Using this idea in the Kingdom's dilemma, one will 
be led to conclude that both players should choose PEACE. Personally, this seems much more appropriate 
to me than the Nash perspective. This alternative viewpoint leads to 
ideas such as superrationality and other decision theories. Classical game theory, however, 
adopts the Nash perspective and studies unilateral deviations and this is the commonly accepted viewpoint.\\

Interestingly enough, the fact that most people who study game theory are prone to accepting the 
Nash perspective is one of the reasons why it is important for us to also adopt the Nash perspective! 
One of the ways I like to convince myself of the Nash perspective is that a Nash equilibrium is 
secure against spies -- if a spy in your kingdom reveals your strategy to the other kingdom 
before the actual simultaneous big "reveal", it won't matter as long as you chose the Nash 
answer since deviation by the other kingdom won't affect you. In general, as long as only one 
kingdom gets this leaked info, Nash equilibria are secure.\\

\noindent A side note is that Nash equilibrium is simply a mathematical definiton and is not prescriptive 
in the sense of "what should players do". Under certain assumptions we can conclude 
that Nash equilibrium is the optimal choice.

\section{Problems}

\begin{enumerate}
\item Derive the mixed strategy equilibrium of Rock--Paper--Scissors.

\end{enumerate}